\section{Identification and the quenched limit}
\label{sec:identification}

The expansion and local comparison now identify the law.  At the
auxiliary scale $q_e$, the first derivative of the $Q_0$ crossing
probability is exactly one at criticality.  Its square-insertion
approximation has a common critical limit, so this identity determines the unknown
site-to-square comparison factor.  We then identify every finite
crossing test, convert from the rectangle to the unit-square
normalization, and pass to the conditional law on $\mathcal K$.
All quantities indexed by $e$ in this section are formed separately
in each model, with the notation of
Sections~\ref{sec:expansion} and~\ref{sec:sewing}.

\subsection{Calibration by one crossing derivative}

Recall that $q_e=(\EE N_e(Q_0))^{-1}$ for the centered auxiliary
rectangle $Q_0$. For its translate centered at a proposed insertion
point, let $S_e$ be the critical single-site kernel and $D_e(h)$ the
critical square-insertion kernel, both with weight
$\delta_1-\delta_0$. These are the nonnegative quantities introduced
in Section~\ref{sec:sewing}. Put
\begin{equation}
 c_e=m_e q_e S_e.
 \label{eq:amplitude-ratio}
\end{equation}
The arm upper bound on $S_e$, together with
$q_e\asymp(m_e a_e)^{-1}$, shows that $c_e$ is bounded.

\begin{lemma}[Common calibration]
\label{lem:common-calibration}
There is a constant $c\in(0,\infty)$ such that $c_e\to c$ in both
models. Moreover, for each fixed degree, assignment of one replica
label to each insertion, and finite signed crossing test, its normalized
all-site kernel integral has a limit common to the two models.
\end{lemma}

\begin{proof}
Fix a compact range of insertion locations separated from one another
and from all polygon boundaries. Write $K_{e,s}(x_1,\ldots,x_k)$ for
the all-site kernel and $K_{e,h}(x_1,\ldots,x_k)$ for the kernel with
all insertions replaced by squares of radius $h$.
Proposition~\ref{prop:replacement}, applied successively at the $k$
locations, gives
\begin{equation}
 (m_e q_e)^k K_{e,s}(\boldsymbol x)
 =c_e^k\frac{K_{e,h}(\boldsymbol x)}{D_e(h)^k}
       +o_h(1),
 \label{eq:all-disk-replacement}
\end{equation}
uniformly on this range, in the ordered sense
$\lim_{h\downarrow0}\limsup_{e\downarrow0}|o_h(1)|=0$.
Here and below a kernel includes its specified signed weights.
To check the factors in this iteration, divide the single replacement
estimate by $D_e(h)\ge c_0 a_{e,h}$. Each already replaced location
contributes at most
$S_e/D_e(h)\le C a_e/a_{e,h}$, while its square kernel costs
$C a_{e,h}$. Thus every accumulated error is $o_h(a_e^k)$ before
normalization. The bound $m_e q_e a_e\le C$ proves
\eqref{eq:all-disk-replacement}, including when $c_e$ is small.

For fixed sufficiently small $h$, Proposition~\ref{prop:critical-tests}
gives common critical limits for $K_{e,h}$ and $D_e(h)$.
The limiting denominator is positive: the square lower bound and the
four-arm lower power bound give $D_e(h)\ge c_h>0$ at this fixed $h$.
These limits hold also for convergent insertion locations and are
continuous on the separated range. Consequently the convergence is
uniform there. Indeed, a failure of uniform convergence would give
a sequence of locations with a convergent subsequence, contradicting
the assertion for convergent locations. Since $\sigma_e$ converges
to Lebesgue measure, integrals against a fixed continuous cutoff
therefore converge too. This applies to the triangular counting
measure as well as to the Poisson location integral.

Apply this observation first with $k=1$, $F=\mathbf1_{C_{Q_0}}$, and
weight $\delta_1-\delta_0$. Its full normalized first coefficient is
\begin{equation}
 m_e q_e\int K_{e,s}(x)\,d\sigma_e(x)
       =q_e\EE N_e(Q_0)=1.
 \label{eq:coefficient-one}
\end{equation}
One may first use the localized box from
Section~\ref{sec:expansion}; its omission contributes $o(1)$ also to
this coefficient. By Lemma~\ref{lem:location-ui}, for every $z>0$
there is a continuous cutoff $0\le\chi\le1$, supported on a compact
range away from the polygon boundary and the box boundary, such that
the integral in \eqref{eq:coefficient-one} changes by at most $z+o(1)$
when multiplied by $\chi$. Nonnegativity is used here.

Let $J_\chi(h)$ be the common limit of
\[
 \int\chi(x)\frac{K_{e,h}(x)}{D_e(h)}\,d\sigma_e(x).
\]
The separate insertion upper bound and the square lower bound imply
$0\le J_\chi(h)\le C_\chi$, uniformly for sufficiently small $h$.
For any subsequential limit $c_*$ of $c_e$,
\begin{equation}
 |1-c_*J_\chi(h)|\le z+o_h(1).
 \label{eq:calibration-cutoff}
\end{equation}
Choose $z<1/4$ and then $h$ small. This proves that every such $c_*$
is bounded below by a positive constant. If $c_+$ and $c_-$ are two
subsequential limits, possibly from different models, the same cutoff
and the same $J_\chi(h)$ work for both. Since all $c_e$ have a common
upper bound $C$, \eqref{eq:calibration-cutoff} also gives
$J_\chi(h)\ge(1-z-o_h(1))/C$. Subtracting the two instances yields
\[
 |c_+-c_-|\le
 \frac{2C(z+o_h(1))}{1-z-o_h(1)}.
\]
First let $h\downarrow0$ and then $z\downarrow0$.
All subsequential limits agree, proving $c_e\to c>0$.

For a general fixed coefficient, restrict its location integral to
a compact separated range, losing at most an arbitrarily small
amount by Lemma~\ref{lem:location-ui}. Equation~\eqref{eq:all-disk-replacement}
and the common critical square limits compare any two subsequential
limits of the restricted coefficient. Their difference tends to zero
as $h\downarrow0$. Removing the cutoff proves uniqueness of the full
coefficient limit across both models. Boundedness follows from
Proposition~\ref{prop:expansion-bound}, so these limits exist.
The factor $1/k!$ in the expansion and the finite sum over replica
assignments do not affect this argument.
\end{proof}

This calibration uses a common macroscopic test, rather than a
comparison of the microscopic scales at equal numerical mesh.
We next convert it to the unit-square normalization of the theorem.

\begin{lemma}[The prescribed normalization]
\label{lem:square-normalization}
There is a common constant $B\in(0,\infty)$ such that, in each model,
\begin{equation}
 q_e\EE N_e(Q_*)\longrightarrow B,
 \qquad r_e=\frac{q_e}{B+o(1)}.
 \label{eq:square-normalization}
\end{equation}
In particular, both prescribed scales tend to zero.
\end{lemma}

\begin{proof}
The first expression is the normalized first coefficient of
$\mathbf1_{C_{Q_*}}$ with weight $\delta_1-\delta_0$.
Lemma~\ref{lem:common-calibration} gives a common finite limit $B$.
For positivity restrict insertion locations to a small open patch
near the center of $Q_*$. The square lower bound of
Lemma~\ref{lem:disk-lower} gives
$K_{e,h}(x)\ge c_1 a_{e,h}$ uniformly on a smaller patch, whereas
$D_e(h)\le C_1a_{e,h}$. Choose a nonzero nonnegative cutoff in that
patch. Equation~\eqref{eq:all-disk-replacement} with $k=1$ and $c_e\to c>0$
shows that this part of the limiting first coefficient is strictly
positive. All omitted parts are nonnegative. Hence $B>0$.
Taking reciprocals gives \eqref{eq:square-normalization};
$q_e\to0$ by Lemma~\ref{lem:pivotal-scale}.
\end{proof}

\subsection{Finite tests and threshold laws}

\begin{proposition}[Common finite-time crossing limits]
\label{prop:nearcritical-tests}
Choose finitely many quads in $\mathcal Q_{\mathrm{rat}}$ and finitely
many real parameters. Their jointly coupled crossing indicators,
with the prescribed scales $r_e$, have a limiting law common to the
two models. The assertion holds for one color replica and for two
independent color replicas conditional on the same geometry.
Every such limiting probability is continuous in the real parameter
vector, and the convergence remains valid for convergent parameter
vectors.
\end{proposition}

\begin{proof}
First use the auxiliary scale $q_e$. Expand the expectation of any
bounded function of the specified indicators as in
Section~\ref{sec:expansion}. The degree-zero term has its common critical limit by
Proposition~\ref{prop:critical-tests}, and every positive-degree
coefficient has a common limit by Lemma~\ref{lem:common-calibration}. The terms with repeated
replica labels at one spatial site tend to zero. The majorant in
Proposition~\ref{prop:expansion-bound} is summable uniformly on every
bounded parameter range, so the whole expectation has a common limit.

At fixed degree the history weights are continuous on their finite
labeled history space, including when times coincide or change order.
The kernel integrals are multilinear in these weights. They therefore
have continuous limiting coefficients, and local uniform summability
makes their sum continuous. The same argument applies to convergent
parameter vectors.

At the prescribed scale, the parameter $\lambda$ corresponds to
$t=\lambda/(B+o(1))$ at the auxiliary scale, by
Lemma~\ref{lem:square-normalization}. The conclusion for convergent
vectors proves the assertion in the required normalization.
\end{proof}

\begin{proof}[Proof of Theorem~\ref{thm:main}]
Finiteness and positivity of the expectations follow from
Lemma~\ref{lem:pivotal-scale}; their inverse scales vanish by
Lemma~\ref{lem:square-normalization}.

At fixed mesh, only finitely many cells meet any particular quad
almost surely. Its crossing indicator is monotone in the parameter,
and the closed mark rule gives, at every finite deterministic $t$,
\begin{equation}
 \{\tau_e(Q)\le t\}=\{C_Q\text{ occurs at }p_e(t)\}
 \quad\text{almost surely}.
 \label{eq:cdf-crossing}
\end{equation}
Thus Proposition~\ref{prop:nearcritical-tests} supplies common limits
of the joint distribution functions of every finite set of thresholds.
These limits are continuous at finite parameter vectors.

For completeness, compactification introduces no ambiguity.
Take any weakly convergent subsequence of threshold laws on $\mathcal K$,
which exists by compactness. At finite continuity points of a limiting
coordinate marginal, \eqref{eq:cdf-crossing} identifies its distribution
function with the continuous limiting crossing function. Squeezing
between such continuity points excludes a finite atom and gives the
same identity at every finite point. Apply this to each coordinate
of a finite projection. Its joint distribution function then agrees
with the one supplied by Proposition~\ref{prop:nearcritical-tests}.
Finite-endpoint rectangles determine measures on
$\overline\RR^k$ as well: their increasing or decreasing limits recover
events involving either endpoint at infinity. All finite projections
are therefore uniquely determined. Since the index set is countable,
they determine the law on $\mathcal K$.

In particular, the triangular laws have a full weak limit, which we
denote by $\nu_\triangle$. The annealed one-replica Voronoi laws have
the same limit. Apply the preceding argument also to two replicas.
For the triangular model their joint law is a product at every mesh,
so its limit is $\nu_\triangle\otimes\nu_\triangle$. Proposition~\ref{prop:nearcritical-tests}
shows that two color replicas on a common Poisson geometry have
this same limit.

Let $f$ be a real continuous function on $\mathcal K$, and put
\[
 Z_{e,f}=\EE[f(\tau_e^V)\mid\eta_e],
 \qquad m_f=\int f\,d\nu_\triangle.
\]
Conditional independence of the two color replicas gives
\[
 \EE Z_{e,f}\longrightarrow m_f,
 \qquad
 \EE Z_{e,f}^{2}
 =\EE[f(\tau_e^{V,1})f(\tau_e^{V,2})]
 \longrightarrow m_f^2.
\]
Hence $Z_{e,f}\to m_f$ in $L^2$ over geometry.

The family
$\mathcal F=\{f:\|f\|_\infty\le1,\ \Lip_{d_{\mathcal K}}(f)\le1\}$
is uniformly bounded and equicontinuous on compact $\mathcal K$.
For any $\zeta>0$ it has a finite uniform $\zeta$-net
$f_1,\ldots,f_m$. For every sampled geometry,
\[
 d_{\mathrm{BL}}\bigl(\Law(\tau_e^V\mid\eta_e),\nu_\triangle\bigr)
 \le 2\zeta+\max_{1\le i\le m}|Z_{e,f_i}-m_{f_i}|.
\]
The maximum tends to zero in probability. Letting $\zeta\downarrow0$
proves \eqref{eq:main-quenched}. All arguments apply along arbitrary
vanishing-mesh sequences, completing the proof.
\end{proof}
